\chapter{Methodology}\label{ch:method}

\section{Data Sources and Preprocessing}\label{sec:data}

\paragraph{Lichess puzzle database.} The open Lichess puzzle export~\cite{lichessdb} (CC0) contains
roughly six million puzzles. Each has a starting position (FEN), a solution in UCI notation, a
Glicko-2 rating with its deviation, a popularity score, a play count and theme tags. A streaming ETL
reads the 1.1\,GB CSV in chunks of one million rows. It removes duplicates, puzzles with rating
deviation above 150, unplayed puzzles and empty solutions, leaving \textbf{5,877,641 puzzles}. At
server start a second quality gate filters on popularity, play count and solution length, and a
bounded, category-stratified sample of 199,998 puzzles is kept in memory.

A \emph{motif-priority} rule maps Lichess's roughly fifty raw theme tags onto the project's 23
categories. Tags arrive in alphabetical order, so an earlier ``first tag wins'' rule let the phase tag
\emph{endgame} outrank real motifs. It mislabelled about 198,000 of 305,000 puzzles then marked
Endgame (Section~\ref{sec:challenges}). The fixed rule resolves ties by specificity: forced mate,
then concrete motif, then generic attack, then endgame type.

\paragraph{Chess.com games.} Games are fetched from the public Chess.com Published-Data
API~\cite{chesscomapi}. The fetcher walks the player's monthly archives from newest to oldest,
caches every response on disk, and waits between requests as the platform's guidance asks.

\paragraph{Application logs.} Every puzzle attempt in the deployed application is stored with its
timestamp, puzzle, category, rating and outcome. For attempts made since the upgrade it also records
the serving policy, the selection probability and the model's prediction before the outcome. At the
time of writing the logs contain 123 attempts by five players.

\paragraph{Research cohort.} Five app users are far too few to evaluate the player models, so a
cohort was collected from the public API. It contains \textbf{300 players}, 60 in each of five rating
bands (below 1000, 1000--1399, 1400--1799, 1800--2199, 2200 and above). Players had to be active (a
rated rapid or blitz game within the last year) and experienced (at least 100 games in that time
control). Candidate players came from three sources:
\begin{enumerate}
  \item \emph{snowball sampling}: opponents of already-collected players whose game rating fell in a
        band that still needed players;
  \item Chess.com's titled-player lists;
  \item random members of the country player lists.
\end{enumerate}
Country lists alone are dominated by beginners, so the upper bands would never fill without the first
two sources. The cost is that the cohort is not a uniform random sample of Chess.com
(Section~\ref{sec:validity}). For each player the newest 60 rated standard-chess games were
analysed, about 18,000 games in total. Players are stored under a salted SHA-256 pseudonym.

\section{Player Modelling and Style Classification}\label{sec:player-model}

\subsection{Move grading by win probability}

Each player move after the first eight half-moves is analysed with Stockfish at a fixed budget of
75,000 nodes, using MultiPV~2 (the best move and the runner-up). Centipawn scores $c$ (from the side
to move) are mapped to win percentages with the curve Lichess uses~\cite{lichessaccuracy}:

\begin{equation}
  W(c) \;=\; 50 + 50\left(\frac{2}{1+e^{-0.00368208\,c}} - 1\right).
\end{equation}

\noindent The loss caused by a move is $L = W(c_{\text{best}}) - W(c_{\text{played}}) \geq 0$. Moves
are graded as inaccuracies, mistakes or blunders at $L\geq 5$, $10$ and $15$ percentage points. Win
probability is used instead of raw centipawns because the same centipawn loss means very different
things in different positions. A 300-centipawn loss is a blunder at $+50$ but only an inaccuracy at
$+900$, where the game is already decided.

\subsection{Opportunities: measuring weakness with a denominator}

A position is \emph{critical}, an opportunity, when the best move beats the runner-up by at least 10
win-percentage points. A simple recapture of a piece the opponent has just taken is excluded, because
it is usually the only move without being a tactic. Each opportunity is labelled with the tactic of
the engine's best line (Section~\ref{sec:features}), and is recorded as a \emph{hit} if the player's
move lost fewer than 5 points.

Counting errors alone is misleading. A player who meets many fork positions will make more fork
errors even if they convert most of them. The quantity of interest is the hit rate per category,
$\text{hits}_k / \text{opportunities}_k$.

\subsection{Weighting and hierarchical shrinkage}

Games are ranked from newest to oldest. Game $g$ receives weight $w_g = 0.5^{\text{rank}/25} \cdot
\tau(\text{time class})$, with $\tau = 0.5$ for bullet, $0.85$ for blitz and $1$ for rapid. Recent
games count more; bullet games count less, because time pressure rather than tactical knowledge drives
many bullet errors.

Each category's weighted hit count is then shrunk towards an \emph{expected} rate, the rate that is
typical for category $k$ at this player's overall level:

\begin{equation}
  e_{u,k} = \sigma\!\bigl(\operatorname{logit}(h_u) + o_k\bigr), \qquad
  \hat r_{u,k} = \frac{\text{hits}_{u,k} + S\,e_{u,k}}{n_{u,k} + S},
\end{equation}

\noindent Here $h_u$ is the player's overall hit rate. The population offset $o_k =
\operatorname{logit}(h_k) - \operatorname{logit}(h)$ says how much harder or easier category $k$ is
for everyone. It is fitted on the research cohort from 23,796 opportunities, for example $+1.71$ for
pawn endgames and $-0.74$ for sacrifices. The strength $S$ is chosen by cross-validation over
players, and was 128--256 in every fold. The original profiler, which shrank towards $h_u$ alone,
misread ``this motif is hard for everyone'' as ``this player is weak at it''.

\subsection{Rating-group classification (Objective O1)}\label{sec:rating-method}

For each cohort player, 18 features are computed from their analysed games, describing \emph{how}
they play:
\begin{itemize}
  \item mean and median win-percentage loss, and mean centipawn loss;
  \item blunders, mistakes and inaccuracies per game;
  \item the share of errors in the opening, middlegame and endgame;
  \item opportunities per game, and the overall hit rate;
  \item hit rates for tactical, mating, endgame and quiet-move groups;
  \item mean game length, win rate, and the share of blitz games.
\end{itemize}
The PGN rating headers of both players are \emph{excluded}, because they would leak the label. The
classifiers are a multinomial logistic regression on standardised features and a random
forest~\cite{breiman2001random}, both from scikit-learn~\cite{pedregosa2011scikit}. They are
evaluated with stratified 5-fold cross-validation over players, repeated ten times, and a
500-shuffle permutation test.

\begin{honest}[Honest status: style]
The portfolio's style classification (positional, attacking, defensive) was not built as a learned
model. A rule-based archetype heuristic (\code{style\_profile.py}) assigns labels such as ``Tactical
Attacker'' from termination patterns and sacrifice rates, but it has no ground truth and was not
validated. Unsupervised style discovery in the manner of Jayasekara~\cite{jayasekara2018classification}
is left as future work.
\end{honest}

\section{Puzzle Representation and Feature Design}\label{sec:features}

A puzzle is a FEN position before the opponent's last move, followed by a solution line in UCI
notation, which alternates the player's moves and forced replies. Its category comes from the
motif-priority rule for Lichess puzzles, or from the tactic tagger for mined puzzles and game
opportunities.

\paragraph{The tactic tagger.} The original labeller looked at the first move only and used a few
geometric tests. It was measured to agree with Lichess themes on only 31\% of puzzles
(Section~\ref{sec:labeller-results}). The replacement tags the \emph{whole line}:
\begin{itemize}
  \item A line the engine scores as a forced mate is a \emph{Mating Pattern}, whatever its first move
        looks like.
  \item A \emph{fork} requires the moved piece to attack at least two targets. Each target must be
        the king, worth more than the attacker, or undefended. The forking piece must not simply be
        lost to a cheaper capture, and a king only counts as a capturer if the square is undefended.
  \item \emph{Pin} and \emph{skewer} come from ray geometry. Walking outward from a line piece, let
        $P_1$ be the first enemy piece and $P_2$ the second. If $P_2$ is the king or worth more than
        $P_1$, it is a pin. If $P_1$ is the king or worth more than $P_2$, it is a skewer.
  \item A \emph{discovered attack} is recorded when another of the player's line pieces newly attacks
        a valuable or undefended target.
  \item \emph{Endgame types} come from the material left on the board (for example, only rooks and
        pawns means a rook endgame), not from which piece moved.
  \item Only the first two player moves may contribute combination motifs; later moves may add only
        promotion or mate. Incidental pins and pawn captures deep in a continuation were the main
        source of false positives.
\end{itemize}
Five categories are never emitted: Attraction, Interference, Clearance, Zugzwang and X-Ray. They need
deeper reasoning about the line, and emitting nothing is preferable to guessing.

\section{PuzzleNet: a Neural Model of Tactic and Difficulty}\label{sec:puzzlenet}

Two measured weaknesses motivate a learned model, and both are documented in
Chapter~\ref{ch:results}. First, every weakness estimate in the system inherits the tactic labeller's
errors, and the hand-written tagger of Section~\ref{sec:features} agrees with Lichess themes on only
58\% of held-out puzzles ($\kappa = 0.52$), while never emitting five of the categories at all. The
label-quality experiment (E3) shows that perfect labels would raise early targeting well beyond what
the tagger achieves, so labelling is a real bottleneck rather than a cosmetic one. Second, a mined
puzzle's difficulty is a closed-form guess from the size of the player's error and the length of the
solution; it is uncalibrated, carries no uncertainty, and 67.5\% of all real attempts were on mined
puzzles.

Both are supervised learning problems for which a large, free, labelled dataset already exists: the
5.88 million Lichess puzzles, each carrying theme tags and a Glicko-2 rating with its own deviation.
\textbf{PuzzleNet} is a single multi-task network trained on that dataset. It reads a position and the
line that follows it, and predicts the line's tactical category, its Lichess theme tags, and how
difficult it is, with an uncertainty attached to that difficulty.

\subsection{Task definition}

Let $s_0$ be a position with the \emph{solver} to move and $L = (m_1, r_1, m_2, r_2, \dots)$ the
principal line from it, alternating solver moves $m_i$ and replies $r_i$. The same pair $(s_0, L)$
occurs in all three places the system needs a label: a Lichess puzzle (the solution line), a critical
position in an analysed game (Stockfish's principal variation), and a mined puzzle (the forced
continuation built by the miner). The network estimates three things:

\begin{align}
  P(\text{category} = c \mid s_0, L), \quad c \in \{1,\dots,24\}, \\
  P(\text{theme}_j = 1 \mid s_0, L), \quad j \in \{1,\dots,66\}, \\
  b \mid s_0, L \;\sim\; \mathcal{N}\!\left(\mu(s_0, L),\; \sigma^2(s_0, L)\right),
\end{align}

\noindent where $b$ is the puzzle's difficulty on the Lichess rating scale. The theme head is an
auxiliary task: theme tags are a richer description of the same line, and predicting them gives the
shared trunk more to learn from than a single label per puzzle.

\subsection{Input representation}

The input is a fixed vector of 3{,}329 features. Its design follows two principles: give the network
the raw board, and also give it the relations a tactician would look for, so that it need not
rediscover the rules of attack and defence from piece positions alone. Whether the second part earns
its place is an ablation, not an assumption (Section~\ref{sec:ablations}).

Everything is encoded \emph{from the solver's point of view}. When Black is to move, squares are
mirrored ($a1 \leftrightarrow a8$) and colours swapped, so that the solver's pieces always occupy the
first six planes and the solver's pawns always advance up the board. A position and its
colour-mirrored twin therefore produce the identical vector, which halves what the network must learn
and is checked by a unit test on real puzzles.

\begin{itemize}
  \item \textbf{Boards} (2{,}304 bits). Piece-square planes ($12 \times 64$) for three positions:
        before the key move, after it, and after the solver's second move.
  \item \textbf{Moves} (864 bits). Six move blocks: the opponent move that created the position, then
        $m_1, r_1, m_2, r_2, m_3$. Each block holds the from-square, the to-square, the moving piece,
        the captured piece, and flags for check, promotion, under-promotion, castling and en passant.
  \item \textbf{Tactical relations} (90 bits). For each of the first three solver moves: which enemy
        piece types the moved piece attacks, whether it attacks two or more valuable targets, whether
        it is left en prise or defended, discovered and double checks, pins it creates and the types
        of pinned pieces, a piece lined up behind an attacked king, queen or rook, whether the move is
        quiet, a retreat, or captures a defender, and whether the position after it is checkmate.
  \item \textbf{Replies} (6 bits) and \textbf{position facts} (43 bits): forced replies and recaptures;
        hanging pieces on both sides; whether the key move is the only check or the only capture;
        castling rights; the material signature of an endgame; a boxed-in king; the number of solver
        moves in the line; and the mate verdict.
  \item \textbf{Continuous features} (27). Material by piece type and the balance before and after the
        line, the number of legal moves, captures and checks available, pressure on the squares around
        the enemy king and its escape squares, and the number of replies the opponent had.
\end{itemize}

No engine evaluation enters the encoding, so all 5.88M puzzles can be encoded in under an hour on a
laptop, and a position can be encoded during game analysis in about three milliseconds. The one
exception is an optional mate flag, which lets game analysis pass Stockfish's mate verdict for a line
that was cut off before the mate appears.

\subsection{Architecture}

A shared trunk of two fully connected ReLU layers ($3{,}329 \to 512 \to 256$) feeds three linear
heads: a 24-way softmax for the category, 66 independent sigmoids for the themes, and two outputs
$(\mu, s)$ for the difficulty, where $\sigma^2 = e^{s}$. The network has 1.86M parameters. Sharing a
trunk across the three tasks is classical multi-task learning~\cite{caruana1997multitask}: the
representation that says \emph{which} tactic a line contains is largely the representation that says
how hard it is to find.

\subsection{Loss: a noise-aware heteroscedastic objective}

For a batch, with $y_c$ the category, $\mathbf{y}_t$ the theme indicators and $y_b$ the observed
rating (standardised), the loss is

\begin{equation}\label{eq:puzzlenet-loss}
  \mathcal{L} = \underbrace{-\log p_{y_c}}_{\text{category}}
  \;+\; \lambda_t \sum_{j} \mathrm{BCE}\bigl(\sigma(z_{t,j}), y_{t,j}\bigr)
  \;+\; \lambda_r \, \underbrace{\mathrm{sg}\!\left(v^{\beta}\right)}_{\text{$\beta$-NLL weight}}
  \cdot \tfrac{1}{2}\left[\log v + \frac{(y_b - \mu)^2}{v}\right],
\end{equation}

\begin{equation}\label{eq:puzzlenet-var}
  v \;=\; \underbrace{e^{s}}_{\text{content uncertainty}} \;+\; \underbrace{\rho^2}_{\text{known rating noise}} .
\end{equation}

\noindent Two choices in this objective matter for how the model is used later.

\paragraph{Known label noise.} A Lichess rating is not the truth about a puzzle; it is a Glicko-2
estimate that comes with its own deviation $\rho$ (a median of 79 rating points, but far larger for
rarely played puzzles). Equation~\eqref{eq:puzzlenet-var} adds that known measurement variance to the
model's own, so that a noisy label pulls the fit less than a precise one, and so that $e^{s}$ is left
to represent only what the content itself leaves uncertain: how hard this puzzle is, given what is on
the board. That is exactly the quantity the recommender needs for a mined puzzle, which has no
Lichess rating at all, and it is why the variance head is not decoration. Section~\ref{sec:integration-irt}
shows where it enters the Bayesian update.

\paragraph{$\beta$-NLL.} Training a Gaussian likelihood with a learned variance has a known failure
mode: the network can reduce the loss on hard examples by inflating their variance, which in turn
scales down their gradient for the mean, so the hardest examples end up teaching the mean the least.
Seitzer et al.~\cite{seitzer2022pitfalls} correct this by weighting each example's loss by a
stop-gradient copy of $v^\beta$. With $\beta = 1$ the mean receives exactly the plain squared-error
gradient; with $\beta = 0$ the plain likelihood is recovered. This work uses the recommended
$\beta = 0.5$, and both extremes are run as ablations.

\subsection{Optimisation and calibration}

The network, its gradients and its optimiser are written from first principles in
NumPy~\cite{harris2020array}; no deep learning framework is used. Two reasons: the deployed
application keeps no heavy dependency and answers in about a millisecond on a CPU, and every step of
the learning algorithm stays visible and testable. Each hand-derived gradient is checked against
central finite differences in the test suite, for every head and for the linear, single-task and
multi-task configurations.

Training uses AdamW~\cite{loshchilov2019decoupled} (decoupled weight decay on weight matrices only),
gradient-norm clipping, and a learning rate that warms up linearly and then decays on a cosine
schedule~\cite{loshchilov2017sgdr}, with batches of 2{,}048 and He initialisation for the ReLU
layers~\cite{he2015delving}. Model selection keeps the checkpoint with the best validation score.

Raw network outputs are then calibrated on the validation split by fitting two scalars: a softmax
temperature $T$ for the category head~\cite{guo2017calibration}, and a variance scale $c$ such that
$v = c\,e^{s} + \rho^2$ for the difficulty head, in the spirit of post-hoc recalibration for
regression~\cite{kuleshov2018accurate}. Both are one-dimensional optimisations of the validation
likelihood, and both are reported before and after in Chapter~\ref{ch:results}.

\subsection{Data, splits and a fair comparison}

Puzzles are split by a hash of their identifier: 2\% validation, 3\% test, the rest for training.
One further set is carved out first. The rule-based tagger was scored in
Section~\ref{sec:labeller-results} on a held-out sample of 30{,}000 uniformly drawn puzzles plus up to
2{,}000 per category; those 59{,}668 puzzles are \emph{excluded from training and model selection
entirely}, so that both labellers can be scored on exactly the same puzzles, neither having seen them.
The tagger's development sample is not excluded, because the stratified draw takes up to 2{,}000 of
each category and the rarest category has only about 2{,}100 puzzles in total; holding out both
samples would have left the network almost no examples of them.

\subsection{Where the network enters the system}\label{sec:integration-irt}

\begin{enumerate}
  \item \textbf{Labelling critical positions.} Game analysis labels the engine's principal variation
        with PuzzleNet instead of the rule-based tagger. Because game analysis sees an engine line
        rather than a puzzle solution, this shift of input distribution is measured separately
        (Section~\ref{sec:nn-deployment}), and the number of plies the network is shown is chosen from
        that measurement.
  \item \textbf{Rating mined puzzles.} A mined puzzle takes $\mu$ as its rating and $\sigma$ as its
        rating deviation, replacing the closed-form heuristic and the blanket deviation of 300 points.
        The deviation is floored at 75 points, the deviation of a well-played Lichess puzzle: a puzzle
        nobody has ever attempted cannot be known more precisely than that, and the model was
        calibrated on Lichess puzzles rather than on puzzles mined from a player's own games.
  \item \textbf{Feeding the Bayesian recommender.} That per-puzzle deviation is the term $v_b$ in the
        IRT update of Section~\ref{sec:rl-setup}: $\kappa = (1 + 3 v_b / \pi^2)^{-1/2}$ scales how far
        one attempt may move a player's ability and category offsets. An uncertain difficulty estimate
        now moves the player's model less, instead of every mined puzzle being treated as equally
        uncertain.
\end{enumerate}

Every one of these paths falls back to the previous behaviour when no model file is installed or when
\code{PUZZLENET=off} is set, and the rule-based tagger remains available through \code{LABELLER=rules}.
The application therefore still runs, unchanged, without the network.

\subsection{What it is compared against}

The category head is compared with the rule-based tagger on the held-out harness, with paired tests on
the same puzzles (McNemar's exact test~\cite{mcnemar1947note} on per-puzzle correctness and bootstrap
intervals for the differences in $\kappa$ and macro recall). The difficulty head is compared with the
miner's formula, with a constant, with a least-squares fit on line length and the mate verdict, and
with gradient-boosted trees~\cite{ke2017lightgbm} trained on the hand-built features alone. Ablations
trained on a common one-million-puzzle subset isolate the contributions of depth, of the hand-built
features, of the raw board, of multi-task learning and of the $\beta$-NLL exponent, and a
learning-curve sweep from 30{,}000 to 5.5 million puzzles shows how much the data volume matters.
Finally, the labels are carried downstream: into the label-quality simulation, into the 300-player
cohort, and into the prequential replay of the real solve logs.


\section{Generative Model Design}\label{sec:generation}

\begin{honest}[Honest status: no generative model was built]
The portfolio planned a generative model fine-tuned with reinforcement learning to create new puzzles,
optimised for uniqueness and counter-intuitiveness, following Feng et al.~\cite{feng2025generating}.
This was \textbf{not implemented}. That approach needs millions of training puzzles, a neural
generator, and an RL loop in which every candidate is scored by engine search, a compute budget far
beyond an individual MSc project. The implemented alternative, \emph{engine-verified mining} from the
player's own games, meets the legality and uniqueness requirements of Objective O2 by construction.
It does not produce puzzles that are new to the world, and it does not optimise
counter-intuitiveness.
\end{honest}

\paragraph{Engine-verified mining.} For each move the player made, four gates must all pass
(the constants are listed in Appendix~\ref{app:constants}):
\begin{enumerate}
  \item \emph{Detection}: the player's move is at least 150 centipawns worse than the engine's best.
  \item \emph{Clarity}: the best move beats the runner-up by at least 100 centipawns, which rules out
        a second solution.
  \item \emph{Forcing}: each opponent reply in the continuation is at least 80 centipawns better than
        its alternative.
  \item \emph{Verification}: a depth-18 re-analysis confirms the solution.
\end{enumerate}
An earlier miner used heuristics without an engine and produced about 40\% false positives; it was
removed. Accepted puzzles are labelled by the tactic tagger. Their difficulty is estimated by a
closed-form heuristic from the size of the error and the length of the solution. That heuristic is
uncalibrated, and replacing it is a stated item of future work (Section~\ref{sec:future}).

\section{Reinforcement Learning Setup}\label{sec:rl-setup}

The recommendation problem is formulated as a multi-armed bandit: one arm per category, with Bernoulli
rewards (solved or not). This is reinforcement learning in its simplest setting, a single state with
repeated actions~\cite{sutton2018reinforcement}. It is \emph{not} deep reinforcement learning: there is
no neural policy, no value network and no multi-step planning.

\subsection{Beta--Bernoulli Thompson Sampling (original policy)}

Each category $k$ has a posterior $\text{Beta}(\alpha_k,\beta_k)$ over its solve rate. The update is
exact by conjugacy: $\alpha_k \leftarrow \alpha_k + y$ and $\beta_k \leftarrow \beta_k + 1 - y$. The
selection rule draws $\tilde\theta_k \sim \text{Beta}(\alpha_k,\beta_k)$ for every category and serves
$\arg\min_k \tilde\theta_k$. This is the \emph{minimum}, not the maximum of standard
TS~\cite{russo2018tutorial}, because the goal is to practise weaknesses. Priors come from the profile,
with a pseudo-count that grows with the evidence: $C_k = \min(10, n_k+2)$, $\alpha_k = 1 + \hat
r_k C_k$, $\beta_k = 1 + (1-\hat r_k) C_k$. An optional discount $\gamma$ pulls every arm back towards
its prior after each update, as in the discounted TS of Raj and Kalyani~\cite{raj2017taming}. It is
switched off by default, for reasons given in Section~\ref{sec:ablations}.

This policy has a structural flaw. Failing a 2200-rated puzzle and failing a 1000-rated one change
$\beta_k$ by the same amount, so categories that happen to be served harder puzzles look weaker.

\subsection{Difficulty-aware IRT learner with Thompson Sampling (new default)}

The new policy keeps a Gaussian posterior $\mathcal{N}(\mathbf{m}, \mathbf{S})$ over $\mathbf{w} =
(\theta, \delta_1,\dots,\delta_{23})$, with full $24\times24$ covariance, for the model in
Equation~\eqref{eq:rasch}. For an attempt in category $k$ at difficulty $b$, whose own uncertainty is
$v_b$, let $\mathbf{x} = \mathbf{e}_\theta + \mathbf{e}_k$. One online Laplace (extended-Kalman) step
gives:

\begin{align}
  \kappa &= \bigl(1 + 3 v_b / \pi^2\bigr)^{-1/2}, \qquad
  p = \sigma\!\bigl(\kappa(\mathbf{x}^\top\mathbf{m} - b)\bigr), \qquad
  h = \kappa^2 p(1-p), \\
  \mathbf{S}' &= \mathbf{S} - \frac{h\,\mathbf{S}\mathbf{x}\mathbf{x}^\top\mathbf{S}}{1 + h\,\mathbf{x}^\top\mathbf{S}\mathbf{x}}, \\
  \mathbf{m}' &= \mathbf{m} + \frac{\kappa\,(y - p)\,\mathbf{S}\mathbf{x}}{1 + h\,\mathbf{x}^\top\mathbf{S}\mathbf{x}}.
\end{align}

\noindent With one parameter this is term for term the Glicko update: $\kappa$ is Glicko's $g(\phi_j)$,
$p$ is $E$ and $h$ is $1/v$. A unit test checks the equivalence to $10^{-9}$ against the Glicko-2
implementation, which itself reproduces Glickman's published worked example (1464.06 / 151.52 /
0.05999)~\cite{glickman2012example}. Because the covariance is full, evidence from a failed puzzle is
split between ``weak overall'' ($\theta$) and ``weak at this motif'' ($\delta_k$). Before each update
the learner adds random-walk drift, $\mathbf{S} \leftarrow \mathbf{S} + \operatorname{diag}(0.02^2,
0.04^2,\dots)$, so that it can follow a player who is improving.

\paragraph{Selection.} Draw $\tilde{\mathbf{w}} \sim \mathcal{N}(\mathbf{m},\mathbf{S})$ and serve
$\arg\min_k \tilde\delta_k$. This is Thompson Sampling on the difficulty-adjusted category offsets.

\paragraph{Priors.} $\delta_k$ is centred on $\operatorname{logit}(\hat r_{u,k}) -
\operatorname{logit}(e_{u,k})$, the player's deviation from what is typical. The information credited
to game evidence is $0.5 \cdot \tfrac{n}{n+S} \cdot n\,\hat r(1-\hat r)$. The factor $0.5$ reflects
that finding a move in a game and solving a flagged puzzle are related but different tasks. The factor
$n/(n+S)$ keeps the prior wide when pooling dominates, so that puzzle attempts can still move it.

\section{Adaptive Curriculum / Recommendation Strategy}

After choosing a category $k$, the learner sets a target difficulty at which the predicted solve
probability equals $p^\star = 0.65$:

\begin{equation}
  b^\star = m_\theta + m_{\delta_k} - \operatorname{logit}(p^\star), \qquad
  r^\star = 1500 + 173.7178\, b^\star ,
\end{equation}

\noindent and the serving window is $r^\star \pm 150$. The choice of $p^\star$ follows the
zone-of-proximal-development argument~\cite{vygotsky1978mind}: hard enough to demand effort, easy enough
to succeed usually. Its sensitivity is tested in Section~\ref{sec:ablations}.

The \emph{serving cascade} looks for an unseen puzzle in this order:
\begin{enumerate}
  \item the player's own mined puzzles in category $k$ and the rating window;
  \item any unseen mined puzzle, 40\% of the time;
  \item the Lichess pool, category $k$, within the window;
  \item the same, with the window widened by $\pm 200$;
  \item any category within the window;
  \item any category at any rating;
  \item a repeat, flagged as such.
\end{enumerate}
Each served puzzle records the policy, the Monte Carlo selection probability of the chosen category,
and the prediction made before the outcome. These make calibration checks and off-policy evaluation
possible later, as in Nie et al.~\cite{nie2026discovering}.

\section{Evaluation Metrics}\label{sec:metrics}

\begin{table}[H]
\centering
\small
\caption{Metrics used in Chapter~\ref{ch:results}.}\label{tab:metrics}
\begin{tabularx}{\linewidth}{p{0.27\linewidth}Y}
\toprule
\textbf{Metric} & \textbf{Definition} \\
\midrule
Strict / lenient agreement & Tagger label equals the Lichess primary category / is among any category the puzzle's themes map to. \\
Cohen's $\kappa$ & Chance-corrected agreement, interpreted on the Landis and Koch scale~\cite{landis1977measurement}. \\
Targeting rate & Share of served categories that are among the player's true three weakest. \\
Regret & $\sum_t (\delta_{k_t} - \min_k \delta_k)$ in logits (Equation~\eqref{eq:regret}). \\
Frustration / productive zone & Share of served puzzles with true $P(\text{solve})<0.30$ / in $[0.50, 0.85]$. \\
Learning gain & Mean improvement of $\delta$ in the initially weak categories (simulation only). \\
Log loss, Brier, AUC & Proper scoring of predicted probabilities on held-out outcomes. \\
Prequential log loss & Each attempt is predicted before the model learns from it~\cite{dawid1984present}. \\
Accuracy, macro-F1 & Rating-band classification; also accuracy within one band. \\
Split-half reliability & Correlation of a player's category deviations between odd and even games. \\
\bottomrule
\end{tabularx}
\end{table}

\noindent\textbf{Statistics.} Policy comparisons use common random numbers, so that every policy meets
the same simulated player for a given seed. They are analysed as paired differences, with paired
$t$-tests, Wilcoxon signed-rank tests, Cohen's $d_z$, bootstrap 95\% confidence intervals and Holm
correction~\cite{holm1979simple} across the policies compared.
